\section{Related Work}
\label{sec:related}

Our analysis sits at the intersection of four lines of research: the
implicit bias and gradient dynamics of cross-entropy training; the
geometry of deep-network loss surfaces; multimodal gradient competition;
and spectral-spatial deep learning for biomedical and remote-sensing
applications. We survey each in turn.

\subsection{Gradient starvation and implicit bias}

Cross-entropy minimization is known to induce structural biases in the
learned representation. Soudry et al.~\cite{soudry2018implicit} proved
that on linearly separable data, gradient descent on logistic loss
converges to the maximum-margin classifier. Gunasekar et
al.~\cite{gunasekar2017implicit} and Lyu and Li~\cite{lyu2020gradient}
extended this picture to deep homogeneous networks. These results
establish that gradient-based training selects \emph{which} solution
among many that fit the training data.

Pezeshki et al.~\cite{pezeshki2021gradient} identified a distinct
phenomenon termed \emph{gradient starvation}: in the Neural Tangent
Kernel regime, dominant features in the input rapidly saturate the
cross-entropy loss, after which weaker but task-relevant features
receive negligible gradient signal. Our Theorem~\ref{thm:twoscale}
adapts this framework to compositional architectures, where the
``dominant features'' are introduced not by the input distribution
but by an architectural capacity asymmetry between modules. The
result is a module-level starvation rather than a feature-level one.

\subsection{Two-timescale stochastic approximation}

The classical machinery for analyzing coupled gradient flows with
disparate timescales is Borkar's stochastic approximation
framework~\cite{borkar2008stochastic}. Heusel et
al.~\cite{heusel2017gans} adapted this to neural network training by
introducing the Two Time-Scale Update Rule for GANs, proving
convergence to a local Nash equilibrium under explicit step-size
disparity. Our setting differs structurally: the timescale separation
in our compositional model is not imposed by the optimizer (the same
learning rate is applied to both modules) but emerges \emph{implicitly}
from the Hessian eigenvalue gap quantified by Theorem~\ref{thm:hessian}.

\subsection{Hessian geometry of deep networks}

The eigenspectrum of the neural-network Hessian has been characterized
empirically and theoretically. Sagun et al.~\cite{sagun2017empirical}
identified a bulk-plus-outliers structure for overparametrized
networks. Pennington and Bahri~\cite{pennington2017geometry} developed
a random matrix theory analysis that predicts spectral densities under
Gaussian assumptions. Karakida et al.~\cite{karakida2019universal}
derived closed-form mean-field expressions for the leading eigenvalue
of the Fisher Information Matrix, showing it scales with the layer
width and parameter count. Theorem~\ref{thm:hessian} applies these
results blockwise to the joint Hessian of a compositional model,
combining them with Schur complement bounds to yield a capacity-ratio
scaling for the joint condition number.

\subsection{Multimodal gradient competition}

The phenomenon most structurally analogous to ours is gradient
competition in multimodal learning. Wang, Tran, and
Feiszli~\cite{wang2020what} showed empirically that joint training of
multimodal classification networks underperforms unimodal training
because different modalities overfit and generalize at different
rates. They proposed Gradient Blending to dynamically reweight
modality-specific losses. Peng et al.~\cite{peng2022balanced}
identified ``modality laziness'' --- a tendency for one modality to
dominate training, leaving the other under-optimized --- and proposed
On-the-fly Gradient Modulation. These works are empirical; our
contribution adapts and extends their framework to the unimodal
\emph{compositional} setting and provides a rigorous geometric
foundation. The recent biomedical application of gradient blending
to multimodal sarcoma prognosis~\cite{bozzo2024multimodal} confirms
the mechanism operates in medical settings, but is restricted to
explicit dual-modality (MRI + clinical) architectures, not the
single-modality compositional pipelines we consider.

\subsection{Spectral-spatial deep learning}

Spectral-spatial deep learning has been an active area for over a
decade, primarily in two communities: biomedical infrared spectroscopy
and hyperspectral remote sensing. Berisha et al.~\cite{berisha2019deep}
established the benefit of combining spectral and spatial features for
FTIR histology and observed that some tissue types could be classified
from spatial features alone, with ``minimal chemical information''
required --- a foreshadowing of the spatial shortcut phenomenon we
diagnose. More recently, O'Leary et
al.~\cite{oleary2026spatial} systematically compared spatial-spectral
architectures for prostate cancer classification in infrared
spectroscopy and reported the striking finding that compressing the
spectral dimension to as few as 16 features had negligible effect on
all models tested. The authors interpreted this as evidence that
tissue classification requires only a limited spectral signature.
Our Theorem~\ref{thm:twoscale} provides an alternative explanation:
end-to-end training fails to extract the spectral signature even when
it is present.

The conceptually nearest work in biomedical imaging is the recent
characterization of shortcut learning in histopathology by Dawood and
Minhas et al.~\cite{dawood2026confounding}, who showed that AI models
predicting molecular biomarkers from H\&E images exploit spurious
correlations with clinicopathological features rather than learning
genuine biological signal. Our paper extends the shortcut-learning
framework into a new modality (spectroscopic imaging) with a different
mechanism (gradient competition driven by architectural asymmetry
rather than confounder leakage).

\subsection{Frozen and random features}

Frozen random features have a long history in machine learning. Saxe
et al.~\cite{saxe2011random} showed that convolutional networks with
random (untrained) filters achieve competitive performance on standard
benchmarks. Rahimi and Recht~\cite{rahimi2007random} established the
theoretical foundations of random features for kernel methods. More
recently, Coil and Cheney~\cite{coil2025freezing} showed that freezing
layers acts as implicit regularization in overparametrized networks.
These works establish that frozen random features are not an oddity
but a principled choice. Our contribution is to identify the specific
mechanism --- gradient starvation under capacity asymmetry --- under
which freezing is not merely competitive but \emph{strictly preferable}
to joint training.
